% Sections 7.1 to 7.5, from lectures/L07/appendix/a_timers.md: one section per part, A.1 to A.5.

\section{Timers: A Counter with a Target}
\label{c7:a1}
A \textbf{timer} answers a different question than the counter from \secref{c6:a1}. A counter
answers ``what is the current count?''; a timer answers ``has a certain amount of time passed?''

It does that with three parts: an internal counter invisible from outside, a target value
\code{TICK_COUNT} fixed per instance, and a comparison that when true clears the counter and pulses
an output flag \code{timeout} for one clock cycle.

Because the system clock ticks at a known, fixed frequency, counting ticks is the same as measuring
elapsed time. On the DE0-CV that is 50,000,000 ticks per second:

\begin{console}
TICK_COUNT = seconds x 50,000,000 - 1
\end{console}

\noindent The \code{- 1} is there because the counter runs \code{0} through \code{TICK_COUNT}
inclusive, so a full period is \code{TICK_COUNT + 1} cycles, not \code{TICK_COUNT}.
Exercise~\ref{c7:ex:3} turns on exactly that, so it is worth reading off this formula rather than
the round number beside it:
\begin{itemize}
  \item \code{TICK_COUNT = 49,999,999}: timeout once per second.
  \item \code{TICK_COUNT = 24,999,999}: twice per second.
  \item \code{TICK_COUNT = 4,999,999}: ten times per second.
\end{itemize}
The rest of this chapter writes the round \code{50,000,000} and \code{5,000,000} where the exact
cycle count does not matter: one tick in fifty million is between 0.02 and 0.2\,ppm at these
values, far below the board oscillator's own tolerance.

This is the inverse of a blocking \code{_delay_ms()}. A blocking delay stops the whole program to
wait; a timer runs permanently in parallel with everything else and simply raises a flag, and
nothing in the design is ever blocked waiting for it.

\section{Building a Timer by Hand in CircuitVerse}
\label{c7:a2}
Before writing any VHDL, build a timer as a gate network, exactly as you did for the combinational
circuits of Chapters~\ref{c1:ch} and~\ref{c2:ch} and the flip-flops of Chapter~\ref{c3:ch}. A timer
is small enough to draw in full, and drawing it is what makes \secref{c7:a3}'s VHDL read as a
description of something you already understand rather than a new idiom to memorize.

Build it at a size you can watch: a \textbf{4-bit counter} with \code{TICK_COUNT = 10}, which is
\code{1010}. The real design counts to 50,000,000; nothing about the structure changes, only the
width.

\textbf{Start from the counter you already built} in \secref{c6:a1}: a 4-bit register \code{Q0} to
\code{Q3} sharing one clock, an adder and a constant \code{1} with the sum fed back to the \code{D}
inputs, and nothing else. If you did not keep it, it is three elements and about five minutes to
redraw. Confirm it still counts and wraps before adding anything, so that when the timer misbehaves
later you know the counter underneath is sound.

Everything here adds \textbf{two things on top of that}: a comparison saying ``the count has
arrived'', and a clear that makes it repeat. That is the entire difference between a counter and a
timer.

\textbf{Add the comparison.} You want \code{timeout} high exactly when the four bits spell
\code{1010}, and the clean way to build ``are these two numbers equal?'' is one \textbf{XNOR gate
per bit}. An XNOR outputs \code{1} when its inputs are the \emph{same}, which is \secref{c1:a1}'s
XNOR column read as an equality test. Wire four, each comparing a counter output against the
corresponding target bit (\code{Q3} against \code{1}, \code{Q2} against \code{0}, \code{Q1} against
\code{1}, \code{Q0} against \code{0}), driving the target side from a CircuitVerse constant rather
than a wire you might later mistake for a signal. Then AND the four outputs together; that single
AND gate is \code{timeout}.

Writing $\odot$ for XNOR, which is the one operator \secref{c1:a2} gave no symbol to:
\[
  \bool{timeout} = (\bool{Q3} \odot 1) \cdot (\bool{Q2} \odot 0) \cdot (\bool{Q1} \odot 1)
    \cdot (\bool{Q0} \odot 0)
\]
All four bits must match at once, which is what ``the counter has reached \code{1010}'' means. This
is a \textbf{general equality comparator}, worth recognizing as its own building block: the same
four XNORs plus an AND compare against any 4-bit value, and only the constants change.

\textbf{Add the clear}, so the timer restarts rather than running on: feed \code{timeout} back to
force the counter to \code{0000} on the next edge, gating each flip-flop's \code{D} input so
\code{timeout = 1} overrides the adder's output. Without it the counter keeps going to \code{1111}
and wraps on its own, so \code{timeout} fires every sixteen edges rather than every eleven. Both
are periodic; only one has the period you asked for.

\textbf{Then confirm, one clock edge at a time:}
\begin{itemize}
  \item The counter advances \code{0000} through \code{1010}, and \code{timeout} is low for every
    count except \code{1010}. Watch the four XNOR outputs as it counts and notice how often three
    are high at once. Three matching bits is not a match; the AND gate is what insists on all four.
  \item On the edge after \code{timeout} fires the counter is back at \code{0000} and
    \code{timeout} is low, so it is high for exactly one clock cycle, never two.
  \item The full cycle is therefore eleven edges, counts \code{0} through \code{10} inclusive,
    repeating indefinitely. That is \secref{c7:a1}'s \code{TICK_COUNT + 1} off-by-one, and this is
    the easiest place in the book to see it, because you can count the edges by eye.
\end{itemize}

\bookimage[htb]{0.6}{lectures/L07/appendix/images/timer_circuit.png}{The comparison logic:
four XNORs against the constants \code{1}, \code{0}, \code{1}, \code{0}, ANDed into
\code{timeout}.}{c7:fig:comparator}

\noindent The drawing in Figure~\ref{c7:fig:comparator} is cropped to the part this section adds.
It labels the counter's bits \code{counter[3]} down to \code{counter[0]} where the text above calls
them \code{Q3} to \code{Q0}, and the increment feeding them is drawn as gate-level carry logic
rather than the single adder block \secref{c6:a1} used, because that is what CircuitVerse gives
you when you build it a bit at a time. The four XNORs, their constants \code{1}, \code{0},
\code{1}, \code{0}, and the AND that combines them are the whole of what is new here.

Two things are worth noticing before the VHDL, because both are what the language buys you:
\begin{itemize}
  \item \textbf{The comparator is hard-wired to one target.} Changing \code{TICK_COUNT} from
    \code{10} to \code{12} means going back into the drawing and flipping two constants. It is a
    small edit, but it is an edit to the \emph{circuit}, and every instance needs its own copy. In
    \secref{c7:a3} the target becomes a generic.
  \item \textbf{The width is hard-wired too.} Counting to 50,000,000 needs 26 flip-flops and a
    26-bit comparator, 26 XNORs across 52 inputs: the same circuit, an unreasonable drawing. In
    \secref{c7:a3} it is the same six lines of VHDL either way.
\end{itemize}
That is the trade this book makes constantly: draw the circuit once, at a size you can see, to
know what the tool is building for you, then let the tool build it at a size you could not have
drawn.

\section{The \code{timer} Module in VHDL}
\label{c7:a3}
The same circuit as a reusable module. The course builds it live, and Exercise~\ref{c7:ex:4} asks
you to write it yourself afterwards, so it is printed here in full rather than shipped as a file:

\begin{vhdlcode}
entity timer is
    generic(TICK_COUNT: natural := 50_000_000);
    port(clock, reset_s2_n, enable: in std_logic;
         timeout                  : out std_logic);
end entity;

architecture behaviour of timer is
-- Internal tick counter.
signal counter: natural range 0 to TICK_COUNT;
begin
    process(clock, reset_s2_n) is
    begin
        if (reset_s2_n = '0') then
            timeout <= '0';
            counter <= 0;
        elsif (rising_edge(clock)) then
            timeout <= '0';
            if (enable = '1') then
                if (TICK_COUNT > counter) then
                    counter <= counter + 1;
                else
                    timeout <= '1';
                    counter <= 0;
                end if;
            end if;
        end if;
    end process;
end architecture;
\end{vhdlcode}

\begin{itemize}
  \item \code{TICK_COUNT} is a \textbf{generic}, not a constant baked into the architecture, so the
    same entity serves any frequency by instantiating it with a different \code{generic map}. The
    internal counter's range is written in terms of it, so the register is sized to the target
    rather than to whatever an unconstrained \code{natural} would give, the same point
    \secref{c6:a1} makes about \code{natural range 0 to 15}. A generic can size a signal precisely
    because it is fixed before synthesis runs.
  \item \code{reset_s2_n} is the already-\emph{synchronized} reset (\secref{c7:a4}), never the raw
    asynchronous \code{reset_n}. A timer never touches a raw asynchronous input directly.
  \item \code{enable} \textbf{freezes} the timer rather than resetting it. While it is \code{'0'},
    \code{counter} holds its value and \code{timeout} stays low, and a disabled timer resumes from
    where it left off rather than restarting. Chapter~\ref{c8:ch}'s \code{fsm_led} depends on this.
  \item \code{timeout} is a single-cycle pulse. The \code{else} branch that sets it is reached on
    only one edge: the edge at which \code{counter} is already sitting at \code{TICK_COUNT} and is
    cleared back to \code{0}. On the following edge the unconditional \code{timeout <= '0';} runs
    with nothing after it to override, so the pulse is exactly one cycle wide.
\end{itemize}
Compare this against \secref{c7:a2}'s gate network. The logic is the same, and \secref{c7:a2}
already named what the generic buys. One difference it did not: \secref{c7:a2}'s \code{timeout} is
a combinational decode, high during the cycle the count reads \code{1010}, while this one is
registered, high one cycle later, during the cycle \code{counter} reads \code{0}. Same period,
different phase, and different glitch behaviour, which is the registered-versus-combinational point
\secref{c8:a5} returns to.

\section{Composing a Full Timer Circuit}
\label{c7:a4}
A timer alone only pulses \code{timeout}; something has to react to that pulse, and a button has to
control the timer. Reusing Chapter~\ref{c4:ch}'s double-flop synchronizer, a complete
button-controlled, timer-driven LED circuit looks like Figure~\ref{c7:fig:floorplan}.

\bookimage{1.0}{lectures/L07/appendix/images/led_toggle_timer.png}{The full circuit at a glance:
synchronized button, timer, and toggled LED, with the detail readable in the two closer views of
Figures~\ref{c7:fig:sync} and~\ref{c7:fig:toggle}.}{c7:fig:floorplan}

\noindent Read that one as a floor plan rather than as wiring. It is a wide circuit, so at page
size its labels are small; what it is for is showing how few pieces there are and roughly where
each one sits. The two closer views below are where the signal names are legible, and between them
they cover everything in it.

Left to right:
\begin{itemize}
  \item \textbf{Synchronization and edge detection} (Figure~\ref{c7:fig:sync}). \code{button_n}
    passes through two flip-flops for metastability protection, then a third lets the circuit
    compare ``now'' against ``previous'' and detect a falling edge, producing a one-cycle
    \code{button_edge_s2} pulse.
  \item \textbf{The timer} (\secref{c7:a3}), enabled or disabled by \code{timer_enabled}, which is
    itself toggled by \code{button_edge_s2}.
  \item \textbf{Toggling the LED} (Figure~\ref{c7:fig:toggle}). The timer's \code{timeout} pulse
    flips the LED's state each time it fires, and the LED is forced off whenever the timer is
    disabled or the circuit is reset.
\end{itemize}

\bookimage{0.9}{lectures/L07/appendix/images/button_sync_and_flank_detection.png}{The synchronizer
and falling-edge detector: \code{button_n} through two flip-flops and a third for the previous
value, ANDed into \code{button_edge_s2}, which toggles \code{timer_enabled}.}{c7:fig:sync}

\bookimage{0.45}{lectures/L07/appendix/images/led_toggle.png}{The LED toggle logic, driven by the
timer's \code{timeout} pulse and forced off while the timer is disabled.}{c7:fig:toggle}

\noindent Two patterns matter here, independent of this circuit:
\begin{itemize}
  \item \textbf{Every asynchronous input is synchronized before it touches any other logic}, and
    every enable/toggle signal derived from a button is a synchronous signal from that point on.
    You reuse this unchanged in \secref{c7:a5} and again for Chapter~\ref{c8:ch}'s state machines.
  \item \textbf{One clock, and everything else is an enable.} \code{timeout} paces the LED by being
    a one-cycle \emph{enable} on a flip-flop that runs on the system clock, not by being wired into
    a clock input. It is tempting to route a counter or timer output into a clock port to ``slow
    things down''; do not. A design with one clock is a design whose timing the tool can analyze,
    and it is the single most common mistake a programmer makes writing a baud generator.
\end{itemize}

\section{FPGA Implementation: A Walking-LED Shift Register}
\label{c7:a5}
\code{walking_led} is the first design that composes three chapters' building blocks into one
hardware-demonstrable circuit: a single lit LED walking along a row, one position per timer tick,
started and stopped with a button. The course builds it live, and Exercise~\ref{c7:ex:6} asks you
to write it yourself afterwards, so it is printed here in full rather than shipped as a file.

It writes almost no new logic, reusing unchanged \code{reset_sync} and \code{button_sync}
(\secref{c4:a5} and \secref{c4:a6}) to synchronize and edge-detect the button, \code{timer}
(\secref{c7:a3}) to pace the rate, and the shift register from \S\S\ref{c6:a2}--\ref{c6:a4} for
the walking itself. That is what the last four chapters have been accumulating: a small set of
modules that compose.

You write all three yourself: \code{reset_sync} and \code{button_sync} in Chapter~4's
Exercises~\ref{c4:ex:8} and~\ref{c4:ex:9}, \code{timer} in Exercise~\ref{c7:ex:4} of this chapter.
Copy your own into the exercise directory before building this. Reusing a module means having one,
and this is the first design that asks you for the ones you already built.

\begin{vhdlcode}
entity walking_led is
    generic(LED_COUNT : natural := 8;
            TICK_COUNT: natural := 25_000_000);
    port(clock, reset_n, button_n: in std_logic;
         led                     : out std_logic_vector(LED_COUNT-1 downto 0));
end entity;

architecture behaviour of walking_led is
signal reset_s2_n, button_edge_s2: std_logic;
signal button_n_v, button_edge_s2_v: std_logic_vector(0 downto 0);
signal shift_enable, shift_tick  : std_logic;
signal shift_reg: std_logic_vector(LED_COUNT-1 downto 0);

begin
    led <= shift_reg;

    button_n_v(0)     <= button_n;
    button_edge_s2 <= button_edge_s2_v(0);

    reset_sync1: entity work.reset_sync
        port map(clock, reset_n, reset_s2_n);

    button_sync1: entity work.button_sync
        generic map(1)
        port map(clock, reset_s2_n, button_n_v, button_edge_s2_v);

    timer1: entity work.timer
        generic map(TICK_COUNT)
        port map(clock, reset_s2_n, shift_enable, shift_tick);

    SHIFT_ENABLE_PROCESS: process(clock, reset_s2_n) is
    begin
        if (reset_s2_n = '0') then
            shift_enable <= '0';
        elsif (rising_edge(clock)) then
            if (button_edge_s2 = '1') then
                shift_enable <= not shift_enable;
            end if;
        end if;
    end process;

    SHIFT_PROCESS: process(clock, reset_s2_n) is
    begin
        if (reset_s2_n = '0') then
            shift_reg <= (0 => '1', others => '0');
        elsif (rising_edge(clock)) then
            if (shift_tick = '1') then
                shift_reg <= shift_reg(LED_COUNT-2 downto 0) & shift_reg(LED_COUNT-1);
            end if;
        end if;
    end process;
end architecture;
\end{vhdlcode}

\noindent Design choices worth noting:
\begin{itemize}
  \item \code{button_n_v}/\code{button_edge_s2_v} bridge a width difference only:
    \code{walking_led}'s button signals are single bits while \code{button_sync}'s ports are
    \code{COUNT}-wide vectors. The two concurrent assignments wire the scalar to element \code{0}
    and back. They cost no hardware.
  \item This is a PISO register (\secref{c6:a4}) in its simplest form: parallel-loaded exactly
    once, on reset, with a single lit bit at position 0, and from then on only shifting. There is
    no runtime \code{load}, because this design never loads a new pattern.
  \item It is a \textbf{circular} shift register. The shift is \secref{c6:a2}'s expression
    unchanged, so the lit bit walks toward the MSB; the only difference from a plain PISO is what
    feeds bit 0, which is the bit falling off the top rather than a \code{serial_in} from
    elsewhere. A one-shot ``shift a pattern out'' register becomes an endlessly repeating pattern
    generator for free, just by recirculating its own output.
  \item \code{shift_enable} is \secref{c7:a4}'s toggle flip-flop reused verbatim: the button's
    synchronized edge pulse flips a single bit, and that bit becomes the timer's \code{enable}.
  \item \code{shift_tick} is the timer's \code{timeout} pulse, reused as the shift register's clock
    enable rather than driving an LED toggle. Same building block, different consumer, and exactly
    the ``one clock, everything else is an enable'' rule from \secref{c7:a4}.
  \item Both \code{LED_COUNT} and \code{TICK_COUNT} are generics, so
    \code{generic map(4, 25_000_000)} gives a 4-LED row and \code{generic map(8, 5_000_000)} a
    faster 100\,ms step, without touching the module.
\end{itemize}
On the board, as the course demonstrates it, the ports are assigned in the Pin Planner
(Appendix~\ref{app:quartus}): \code{clock} to the 50\,MHz oscillator, \code{reset_n} and
\code{button_n} to two push buttons, and \code{led} to a row of onboard LEDs. Press the button once
to start the bit walking, again to freeze it.

You can watch the same behaviour without a board by running the reference testbench, which
overrides both generics so a full lap takes a handful of clock cycles:

\begin{shell}
cd lectures/L07/exercises/walking_led
cp ../../../L04/exercises/reset_sync/reset_sync.vhd .   # the three you wrote yourself
cp ../../../L04/exercises/button_sync/button_sync.vhd .
cp ../timer/timer.vhd .
ghdl -a --std=93 reset_sync.vhd button_sync.vhd timer.vhd \
                 walking_led.vhd walking_led_tb.vhd
ghdl -e --std=93 walking_led_tb
ghdl -r --std=93 walking_led_tb --assert-level=error --stop-time=10ms
\end{shell}
